\section{Conclusions}
\label{sec:conclusions}

\subsection{Summary of Contributions}
\label{subsec:contributions}

This report addressed a documented, economically significant problem in the building industry:
manual BACnet point mapping costs 0.5--1.5\% of construction value per building, blocks software
interoperability, and is identified in peer-reviewed literature and DOE policy reports as a major
impediment to building automation innovation.
The same problem directly prevents the large-scale enrollment of commercial buildings into Virtual
Power Plant and Non-Wire Alternative programs, because each building requires a custom, per-project
integration effort before it can participate in grid services.
No widely-deployed automated solution for this problem existed prior to this work.

SI-Mapper demonstrates that agentic AI can automate the creation of ASHRAE 223P semantic building
models from raw inputs available on any real building project: HVAC engineering drawings and
BACnet point exports.
The system combines multimodal drawing recognition, using foundation models capable of interpreting
HVAC schematic symbols directly from images, with automated BACnet point-to-equipment mapping and
formal ASHRAE 223P ontology generation in a single, end-to-end pipeline.
The agent validates its own ontology output through iterative self-correction, looping until the generated
Turtle (TTL) file executes without errors and the resulting graph loads cleanly into the Neo4j
semantic store.
This pipeline replaces a workflow that previously required days of expert HVAC engineering labor.
Two experiments on real air handling unit installations confirm the pipeline's end-to-end
viability: System 1 (21 components, 44 BACnet points, 295 nodes and 949 relationships in the
resulting graph) was completed in 17 minutes 11 seconds at an estimated cost of \$11.30 USD;
System 3 (27 components, 34 BACnet points, more complex heat recovery loop) was completed in
12 minutes 24 seconds at \$10.85 USD.
In both runs the agent identified and corrected ontology code errors autonomously, three
iterations for System 1, one for System 3, without requiring human intervention.
Both sessions used Gemini 3.1 Pro Preview and stayed well within the model's context window
throughout.

The development process itself is a contribution.
Nine experimental iterations, from simple LangChain function-calling in 2025 to the final
skills-based architecture on Gemini 3.1 and Claude Sonnet 4.6 in 2026, constitute a systematic
record of how agentic AI architectures for domain-specific technical tasks evolved as foundation
model capabilities improved.
The progression reveals a counterintuitive architectural lesson: simpler architectures became
more viable, not less, as models improved.
Multi-agent designs with dedicated Planner, Actor, and Reviewer sub-agents, initially necessary
to keep the pipeline on track, became unnecessary once models reached sufficient instruction-following
capability.
The final skills-based single master agent outperformed every multi-agent architecture tried
while being dramatically simpler to maintain, debug, and extend.

The technology stack assembled for SI-Mapper, Google ADK, CopilotKit, Neo4j with Neosemantics
(n10s), and Sigma.js, provides a reusable reference architecture for agentic AI applications
requiring real-time operator observability, semantic graph storage, and interactive graph
visualization.
The combination of ADK for agent orchestration, CopilotKit for the Agent-to-UI streaming protocol,
and Neo4j for native RDF graph storage is not an obvious combination; the design decisions that
led to it are documented in Section~\ref{sec:implementation} as a practical guide for future
implementors.

\subsection{Limitations}
\label{subsec:limitations}

\textbf{Model dependency.}
The system's performance is directly coupled to the capabilities and availability of specific
foundation models, Gemini 3.1 and Claude Sonnet 4.6 as of this writing.
Model API changes, pricing changes, or capability regressions in future model versions could
affect pipeline reliability without any change to the SI-Mapper codebase.
The skills architecture mitigates this risk somewhat by encapsulating domain knowledge in
plain-text instruction sets that can be updated independently of model versions, but the
underlying dependency on external model providers remains.

\textbf{Limited validation scope.}
The current experimental scope covers two real HVAC systems (System 1 and System 3), both air
handling units drawn from the same project context and following the same Quebec French BACnet
naming conventions.
Both experiments succeeded, the agent produced a valid ASHRAE 223P Turtle file in a single
session in each case, but two systems from a single contractor and region are insufficient to
establish statistical confidence in accuracy or generalizability.
No failed runs were observed, though minor ontology code errors (incorrect imports, missing
constructor parameters, operator syntax) did occur and were autonomously corrected in all cases.
Whether such errors would be recoverable on systems with significantly different drawing styles,
denser point lists, or more unusual equipment topologies is not yet known.

\textbf{Domain specificity.}
The current skills and tool set are tailored to HVAC air handling units with BACnet points
in French-language naming conventions, reflecting the Quebec installation context where this
work was developed.
Extending SI-Mapper to other building system types (VAV boxes, chiller plants, boiler plants),
other control protocols (Modbus, KNX, BACnet/IP vs.\ BACnet MSTP), or BACnet point naming
conventions from other regions and contractors would require additional skill development and
validation.
The architecture supports this extension naturally, new skills are plain-text instruction
sets, but the domain knowledge required to write accurate skills is non-trivial.

\textbf{ASHRAE 223P maturity.}
The standard is currently in its second advisory public review and is not yet ratified.
The ontology generation logic in SI-Mapper is calibrated against the current advisory draft.
Changes to the standard during ratification, particularly to connection point types, medium
definitions, or telemetry binding syntax, could require updates to the ontology generation
skills and the underlying ASHRAE 223P Python library used to produce the TTL output.

\subsection{Future Work}
\label{subsec:future-work}

\textbf{Production validation at scale.}
The most important next step is continue developing SI-Mapper and test the model against a larger corpus of real buildings,
a target of 20--50 HVAC systems spanning multiple system types, contractors, and naming conventions.
This would establish statistical confidence in accuracy and cost savings, identify edge cases
not encountered in the current development set, and produce the benchmark data needed for
peer-reviewed publication.

\textbf{Multi-system support.}
The current skills cover air handling units.
A production-viable system would need skills for VAV terminal boxes, chiller plants, boiler
plants, heat pump systems, and mechanical ventilation units.
Each system type has its own drawing conventions, point naming patterns, and ASHRAE 223P
topology rules.
The skills architecture makes this extension tractable: each system type requires a new
skill file rather than a new agent or pipeline redesign.

\textbf{Automated quality metrics.}
The current evaluation of generated ASHRAE 223P models is qualitative.
A production evaluation framework would include automated comparison tools: graph diffing
against a human-authored ground-truth model, quantitative scores for completeness (fraction
of expected equipment nodes present), accuracy (fraction of BACnet references correctly
assigned), and topology correctness (fraction of directed connections matching the source
drawing).
These metrics would make it possible to track regression across model updates and skill
revisions without manual inspection.

\textbf{Multi-language support.}
Adapting the BACnet point mapping skills to recognize English, Spanish, and other naming
conventions beyond French would expand SI-Mapper's applicable geography significantly.
The mapping skills are the primary location where naming convention knowledge is encoded;
adding language variants is a skill-writing task rather than an architectural change.

\textbf{Real-building demonstration pilot.}
The highest-impact near-term validation step is to connect a generated ASHRAE 223P model
to live BACnet points on a real building and operate against real data.
The pilot would proceed in two stages: first, bind the ontology nodes to their physical
counterparts and verify that the semantic graph can be used to query live sensor readings,
confirming that the ASHRAE 223P ontology interface is sufficient for building monitoring; second,
issue simple supervisory control commands through the same interface.
This would validate the central hypothesis of SI-Mapper: that the generated semantic model
is a usable operational interface to monitor and control buildings, not merely an enhanced documentation artifact.